\section{Methodology}
\label{sec:methodology}

\subsection{The StabilitySection Framework}

StabilitySection is not a new redistricting algorithm.
It is an analysis framework that runs GeoSection once per census year for
each state and compares outputs along four dimensions:
natural ratio, bisection tree structure, district assignments, and partisan seats.

\begin{algorithm}
\caption{StabilitySection}
\label{alg:stabilitysection}
\begin{algorithmic}[1]
\Require State $s$, census years $Y \subseteq \{2000, 2010, 2020\}$,
         seat count function $k(\cdot)$, seed count $N$
\Ensure StabilityReport for state $s$
\For{$y \in Y$}
  \State $G_y \gets \textsc{LoadGraph}(s, y)$ \Comment{tract adjacency graph}
  \State $\mathcal{M}_y \gets \textsc{GeoSection}(G_y, k(s, y), N)$
  \State $r_y \gets \mathcal{M}_y.\texttt{level1\_ratio}$
  \State $T_y \gets \mathcal{M}_y.\texttt{bisection\_tree}$
  \State $A_y \gets \mathcal{M}_y.\texttt{tract\_assignments}$
  \State $p^*_y \gets \mathcal{M}_y.\texttt{lorenz\_pstar}$
  \State $\text{seats}_y \gets \textsc{PartisanEval}(A_y)$
\EndFor
\State \Return $\textsc{ComputeCSS}(\{r_y\}, \{T_y\}, \{A_y\}, \{p^*_y\}, \{\text{seats}_y\})$
\end{algorithmic}
\end{algorithm}

\subsection{The Census Stability Score}
\label{sec:css-definition}

\begin{definition}[Census Stability Score]
For a state $s$ with GeoSection outputs across years $Y$, the
Census Stability Score is:
\[
  \mathrm{CSS}(s) = 0.5 \cdot s_{\text{seat}}(s)
                  + 0.3 \cdot s_{\text{ratio}}(s)
                  + 0.2 \cdot s_{\text{gap}}(s),
\]
where each component is a normalised stability score in $[0,1]$.
\end{definition}

The weights reflect the legal importance ordering.
Partisan seat counts ($0.5$ weight) are the primary legal question: does the
algorithm produce the same D/R outcome regardless of which census year's data
is used?
Natural ratio stability ($0.3$ weight) measures geographic structure
stability: does the algorithm identify the same first-level split?
Gap magnitude similarity ($0.2$ weight) measures the consistency of the
proportionality deficit, relevant for establishing that any gap is a durable
geographic feature, not a census-specific artefact.

\paragraph{Seat stability.}
\begin{equation}
  s_{\text{seat}}(s) = \mathbf{1}\bigl[\text{dem\_seats}_y = \text{dem\_seats}_{y'}
                       \text{ for all } y, y' \in Y\bigr]
\end{equation}
A binary indicator: 1 if the Democratic seat count is identical across all
census years, 0 if it differs by even one seat.
For states where the seat count $k$ changes between censuses (e.g.\ TX grew from
30 to 38), $s_{\text{seat}}$ is computed on the seat \emph{share}
$\text{dem\_seats}/k$ rather than raw count, with a tolerance of $\pm 0.04$
($\approx \pm0.5$ seats for $k=12$).

\paragraph{Ratio stability.}
\begin{equation}
  s_{\text{ratio}}(s) = \mathbf{1}\bigl[r_y = r_{y'}
                       \text{ for all } y, y' \in Y\bigr]
\end{equation}
A binary indicator: 1 if the GeoSection first-level natural ratio is identical
across all census years.
For states where $k$ changes, we compare the \emph{normalised ratio}
$\min(r_y^L, r_y^R) / k_y$ rather than raw integers.

\paragraph{Gap similarity.}
\begin{equation}
  s_{\text{gap}}(s) = 1 - \frac{\max_{y,y' \in Y}
                         |\mathrm{gap}_{y} - \mathrm{gap}_{y'}|}{100}
\end{equation}
where $\mathrm{gap}_y$ is the proportionality gap in percentage points in year
$y$.
A gap that changes by 5\,pp between census years contributes
$s_{\text{gap}} = 0.95$.
A gap that changes by 20\,pp contributes $s_{\text{gap}} = 0.80$.

\paragraph{CSS thresholds.}
\begin{center}
\begin{tabular}{lll}
\toprule
CSS range & Category & Legal interpretation \\
\midrule
$\geq 0.90$ & Highly stable & Strongest geographic determinism claim \\
$0.70$--$0.90$ & Moderately stable & Seat-stable; some ratio or gap drift \\
$0.50$--$0.70$ & Conditionally stable & Seat changes in some census years \\
$< 0.50$ & Structurally volatile & Different outcomes across censuses \\
\bottomrule
\end{tabular}
\end{center}

\subsection{District Jaccard Similarity}
\label{sec:jaccard}

For states with identical $k$ across census years, we additionally compute
the district assignment Jaccard similarity between years.
This measures how much the actual tract-to-district assignments shift, even
when the seat count and natural ratio remain the same.

\begin{definition}[District Jaccard Similarity]
Given tract assignments $A_y$ and $A_{y'}$ for two census years, with
$k$ districts in both years:
\[
  J(A_y, A_{y'}) = \frac{\sum_{d=1}^{k} \text{pop}(d_y \cap d'_{\sigma(d)})}
                        {\sum_{d=1}^{k} \text{pop}(d_y \cup d'_{\sigma(d)})},
\]
where $\sigma(d) = \arg\max_{d''} \text{pop}(d_y \cap d''_{y'})$ maps each
year-$y$ district to the best-matching year-$y'$ district by population
overlap.
\end{definition}

Note that Jaccard similarity requires a common unit of comparison.
Census tracts are not spatially identical across years (Type II change).
We address this by mapping 2000 and 2010 assignments to 2020 tracts via
spatial intersection weighted by population, treating 2020 tracts as the
reference geography.

\subsection{Lorenz Drift as a Stability Proxy}
\label{sec:lorenz}

GeoSection's natural ratio is determined by the \emph{optimal point} $p^*$
on the population-geography Lorenz curve: the point where the marginal
reduction in normalised edge-cut from increasing the ratio is zero.
A state's Lorenz curve encodes how clustered its population is relative to
its geographic area.

\begin{definition}[Lorenz Drift]
For a state $s$ and census years $y, y'$:
\[
  \Delta p^*(s, y, y') = |p^*_y - p^*_{y'}|
\]
where $p^*_y$ is the population fraction corresponding to the winning ratio
side at census year $y$.
For a $1:k-1$ split in a $k$-district state, $p^* = 1/k$.
For a $j:(k-j)$ split, $p^* = j/k$.
\end{definition}

\begin{proposition}[Lorenz Drift Predicts CSS]
\label{prop:lorenz}
For a state with isoperimetric gap $\Delta \mathrm{EC}_{\text{norm}} > \epsilon$
between the winning and second-best ratio at year $y$, the ratio is stable
at year $y'$ whenever:
\[
  \Delta p^*(s, y, y') < \frac{\epsilon}{2 \cdot C_{\text{pop}}},
\]
where $C_{\text{pop}}$ is the maximum rate of change of normalised edge-cut
per unit of population fraction shift.
\end{proposition}

This proposition formalises the intuition that states with large isoperimetric
gaps (clear winners) are more robust to population perturbation than states
with narrow margins.
We use $\Delta p^*$ as a computationally accessible proxy for CSS when
full cross-census runs are unavailable: states where $p^*$ shifts by more
than $0.05$ between available census years are flagged as potentially volatile.

\subsection{Stability Decomposition: Type I vs.\ Type II Change}
\label{sec:decomposition}

To separate population-shift effects (Type I) from tract-redesign effects
(Type II), we run a third experiment for a 10-state subset:

\begin{enumerate}
\item \textbf{Full cross-year stability}: GeoSection on the actual 2000 graph
  with 2000 populations, compared to the actual 2020 graph with 2020
  populations.
  This is the standard CSS measurement.

\item \textbf{Population-only perturbation}: GeoSection on the \emph{2020}
  graph with 2000 populations (census population totals re-allocated to
  2020 tracts by spatial intersection).
  The graph topology is fixed; only the node weights change.
  The CSS difference between (1) and (2) is the contribution of tract
  boundary redesign.
\end{enumerate}

For most states, we expect the population-only perturbation to show higher
stability than the full cross-year comparison, because fixing the graph
topology removes the contribution of tract splits and merges.
The decomposition is of independent scientific interest: states where
tract-redesign contributes significantly to instability are states where
the Census Bureau's geographic decisions (not just population growth) affect
the algorithmic output.

\subsection{CSS for Two Census Years}
\label{sec:two-year-css}

With complete 2020 data and partial 2000 data, we compute a two-census CSS:
\[
  \mathrm{CSS}_{2\text{yr}}(s) = 0.5 \cdot s_{\text{seat}}^{(2000,2020)}
                                + 0.3 \cdot s_{\text{ratio}}^{(2000,2020)}
                                + 0.2 \cdot s_{\text{gap}}^{(2000,2020)}.
\]
This covers a twenty-year span.
The full three-census CSS will incorporate 2010 when that sweep completes:
\[
  \mathrm{CSS}_{3\text{yr}}(s) = 0.5 \cdot \bar{s}_{\text{seat}}
                                + 0.3 \cdot \bar{s}_{\text{ratio}}
                                + 0.2 \cdot \bar{s}_{\text{gap}},
\]
where each component is the average across all three pairwise year comparisons.

\paragraph{Handling seat-count changes.}
Fourteen states changed their congressional delegation size between 2000 and
2020 (gained or lost seats due to apportionment).
For these states, seat stability is measured on seat-share with tolerance
$\pm 1/k_{\max}$, and ratio stability is measured on the normalised ratio
as described above.
States that both changed seat count and changed natural ratio are assigned
$s_{\text{ratio}} = 0$ only if the normalised ratio change exceeds the
seat-change threshold; if the ratio shifted predictably with the new seat
count, partial credit is given.
